\textbf{(CAP-COLUNI 2018 -- Q14)} Uma quadra de futebol foi construída com as seguintes dimensões: 10 metros de largura e 20 metros de comprimento, conforme a figura abaixo. Obteve-se por círculo central (região circular no centro da quadra), um círculo $C$ de raio 2 metros e, por áreas (regiões nas quais os goleiros podem tocar a bola com as mãos), dois semicírculos, $S_1$ e $S_2$, de raio 4 metros.

Essa quadra será reformada, aumentando-se, proporcionalmente, seu comprimento e sua largura, mas mantendo-se inalterados os raios de $C$, $S_1$ e $S_2$, de tal modo que o percentual ocupado pela soma das áreas de $C$, $S_1$ e $S_2$ seja igual a 10\% da área da nova quadra.

Usando $\pi=3$ e $\sqrt3=1{,}7$, é CORRETO afirmar que a diferença entre o novo comprimento e a nova largura dessa quadra será, em metros, de:

\includegraphics{./figs/2018-Q14.png}

\begin{enumerate}
	\item 11
	\item 13
	\item 15
	\item 17
\end{enumerate}
